\section{Countable-Model Spectra and Infinitary Isomorphism}
\label{sec:p03-n-a-004-02}

The number of isomorphism types is not the number of labeled copies.
Counting countable models therefore involves both a coding space
and an equivalence relation on it. Morley's 1970 threshold concerns
more than \(\aleph_1\) isomorphism types, not merely an uncountable
set of codes \cite{Morley1970CountableModels}.
The construction below supplies the infinitary isomorphism and
coding interfaces used in such spectrum arguments.

\subsection{Codes and the threshold}
\label{sec:p03-morley-codes}

\begin{hbdefinition}{def:p03-morley-coding}{The fixed-domain coding space}
Let \(L\) be an at most countable first-order language of
finite-arity relations, functions and constants, with logical
identity. Structures are nonempty. Write \(X_L\) for its
structures with domain \(\omega\), coded as
\[
 X_L=
 \prod_{R\in L}2^{\,\omega^{\operatorname{arity}(R)}}
 \ \times\
 \prod_{f\in L}\omega^{\,\omega^{\operatorname{arity}(f)}}
 \ \times\
 \prod_{c\in L}\omega .
\]
Positive-arity functions are in the second product;
individual constants are in the third. Nullary relations have
one Boolean coordinate. Each discrete coordinate has its
discrete topology and the whole space has the product topology.
Codes record all interpretations; logical identity is actual equality.
Let \(X_{L,m}\) be the analogous space on
\(\{0,\ldots,m-1\}\), for \(m\geq1\).
The topological sum \(X_L^{\leq\omega}\) of these spaces and
\(X_L\) includes all nonempty at most countable structures.

For a class \(K\subseteq X_L\), put
\(I(K)=|K/{\cong}|\). For a sentence \(\sigma\) or theory \(T\),
\(K_\sigma\) or \(K_T\) is its model class on \(\omega\).
These countably infinite-domain counts are kept distinct
from the all-at-most-countable convention of
\cref{def:p03-vaught-spectrum-notation}.
\end{hbdefinition}

\begin{hblemma}{thm:p03-morley-code-space}{proved=annex:p3.morley.codes}{The coding spaces are Polish and have the continuum upper bound}
The spaces \(X_L\), \(X_{L,m}\) and \(X_L^{\leq\omega}\)
are Polish: separable and completely metrizable.
Their cardinalities are at most \(2^{\aleph_0}\).
Consequently any class of at most countable \(L\)-structures
has at most \(2^{\aleph_0}\) isomorphism types.
For an infinite language, even a fixed finite domain can
have continuum many nonisomorphic structures; finite size
alone is not a countable bound on the number of theories.
\end{hblemma}

\begin{hbtheorem}{thm:p03-morley-firstorder-threshold}{cited=Morley1970CountableModels}{Morley's first-order threshold}
For a complete satisfiable first-order theory \(T\) in a
countable language with infinite models,
\[
 I(K_T)>\aleph_1
       \quad\Longrightarrow\quad I(K_T)=2^{\aleph_0}.
\]
This is a first-order instance of the 1970 threshold.
It leaves the case \(I(K_T)=\aleph_1<2^{\aleph_0}\)
untouched; it does not replace \(>\aleph_1\) by
\(>\aleph_0\).
\end{hbtheorem}

Morley's published introduction also announces extensions to
\(L_{\omega_1,\omega}\) and pseudo-axiomatic classes in that
logic \cite[p.~14]{Morley1970CountableModels}.
An arbitrary collection of infinitary sentences must not
be treated as one countable conjunction without a countability
condition. The precise formulation matters because there are
uncountably many infinitary sentences.
The first-order cardinal conjecture, its perfect-set
strengthening and the dated \(\omega\)-Vaught comparison
are distinguished in
\cref{def:p03-vaught-variants,def:p03-vaught-scott-comparison}.
The edition's selected status is in
\cref{sec:p03-vaught-status-2026}; no different rank convention
or later version is substituted here.

\subsection{Infinitary truth in the coding space}
\label{sec:p03-morley-borel-truth}

\begin{hbdefinition}{def:p03-morley-infinitary}{Countable connectives and formula rank}
The logic \(L_{\omega_1,\omega}\) permits countable
conjunctions and disjunctions, each with only finitely
many free variables in total, and finite strings of
individual quantifiers. Formula trees are well founded.
Define the auxiliary quantifier rank by
\[
 \operatorname{qr}(\theta)=0\ \text{for atomic }\theta,\quad
 \operatorname{qr}(\neg\varphi)=\operatorname{qr}(\varphi),
 \quad
 \operatorname{qr}\!\left(\bigwedge_i\varphi_i\right)
              =\sup_i\operatorname{qr}(\varphi_i),
 \quad
 \operatorname{qr}(\exists y\,\varphi)
              =\operatorname{qr}(\varphi)+1.
\]
Disjunction has the same rank clause as conjunction.
All these ranks are countable ordinals.
A set is Borel if it belongs to the sigma-algebra
generated by the open sets.
\end{hbdefinition}

\begin{hblemma}{thm:p03-morley-borel-truth}{proved=annex:p3.morley.borel-truth}{Each fixed infinitary formula defines a Borel truth set}
For a fixed \(L_{\omega_1,\omega}\)-formula
\(\varphi(\bar x)\) and each fixed tuple
\(\bar a\in\omega^n\),
\[
 \{M\in X_L:M\models\varphi(\bar a)\}
\]
is Borel. The satisfaction relation is also Borel in
\(X_L\times\omega^n\), with the tuple space discrete.
The same holds in each finite-domain component.
Every countable collection of infinitary sentences can
be conjoined into one such sentence, including the
empty collection with conjunction \(\top\).
Its fixed-domain model class is Borel.
\end{hblemma}

\subsection{Transfinite back-and-forth and the Scott construction}
\label{sec:p03-morley-scott}

\begin{hbdefinition}{def:p03-morley-approximants}{Pointed back-and-forth approximants}
For equal-length finite tuples in structures \(M,N\), define
\((M,\bar a)\equiv_\alpha(N,\bar b)\), for \(\alpha<\omega_1\):
\begin{enumerate}
\item At \(0\), the tuples agree on every atomic formula
      in their slots.
\item At \(\alpha+1\), they are \(\equiv_\alpha\)-equivalent,
      and every \(c\in M\) has some \(d\in N\) with
      \((M,\bar a c)\equiv_\alpha(N,\bar b d)\),
      while every \(d\in N\) has some such \(c\in M\).
\item At a nonzero limit \(\lambda\), all
      \(\equiv_\beta\), \(\beta<\lambda\), hold.
\end{enumerate}
Tuples may repeat; atomic equality identifies the same
repetitions on each side. Empty tuples are included.
Write \(E_\alpha(M,N)\) for the empty-tuple relation;
its Borel claims below concern codes on \(\omega\).
This auxiliary ordinal analysis is not an identification
with the parameterized or unparameterized Scott-rank
conventions in \cref{def:p03-vaught-scott-comparison}.
\end{hbdefinition}

\begin{hblemma}{thm:p03-morley-approximant-properties}{proved=annex:p3.morley.approximants}{Equivalence, Borelness and preservation of ranked formulas}
The pointed relations are equivalence relations in each
tuple arity and decrease with \(\alpha\).
For coded structures on \(\omega\), each fixed
countable-ordinal relation is Borel.
If \((M,\bar a)\equiv_\alpha(N,\bar b)\), the tuples
agree on every infinitary formula of quantifier rank
at most \(\alpha\).
\end{hblemma}

\begin{hbdefinition}{def:p03-morley-scott-formulas}{The formulas attached to one countable structure}
Fix an at most countable \(M\) and a finite tuple \(\bar a\).
Let \(\theta_{\bar a}^{0}(\bar x)\) be the conjunction
of the atomic formulas true at \(\bar a\) and the
negations of those false there. Define
\[
\begin{split}
 \theta_{\bar a}^{\alpha+1}(\bar x)
   ={}&\theta_{\bar a}^{\alpha}(\bar x)
       \land\bigwedge_{c\in M}\exists y\,
                         \theta_{\bar a c}^{\alpha}(\bar x,y)\\
      &\land\forall y\bigvee_{c\in M}
                         \theta_{\bar a c}^{\alpha}(\bar x,y),\\
 \theta_{\bar a}^{\lambda}(\bar x)
   ={}&\bigwedge_{\beta<\lambda}
                         \theta_{\bar a}^{\beta}(\bar x).
\end{split}
\]
All bound variables are renamed to avoid capture.
Each indexing set is countable and each formula has
only its displayed finite slots free.
\end{hbdefinition}

\begin{hblemma}{thm:p03-morley-scott-stabilization}{proved=annex:p03-morley-scott-stabilization}{Semantic representation and a countable stabilization stage}
For every structure \(N\),
\[
 N\models\theta_{\bar a}^{\alpha}(\bar b)
       \quad\Longleftrightarrow\quad
 (M,\bar a)\equiv_\alpha(N,\bar b).
\]
There is a countable ordinal \(\delta_M\) at which
all finite-tuple approximants \emph{within \(M\)}
stabilize:
\[
 (M,\bar a)\equiv_{\delta_M}(M,\bar b)
       \quad\Longrightarrow\quad
 (M,\bar a)\equiv_{\delta_M+1}(M,\bar b)
\]
for all equal-length tuples.
This countable stage depends on \(M\); no uniform
countable bound for all structures is asserted.
\end{hblemma}

\begin{hbtheorem}{thm:p03-morley-scott-sentence}{proved=annex:p3.morley.scott-sentence}{A full countable-isomorphism sentence}
The \(L_{\omega_1,\omega}\)-sentence
\[
 \sigma_M=
 \theta_\varnothing^{\delta_M}
 \land\bigwedge_{n<\omega}\ \bigwedge_{\bar a\in M^n}
   \forall\bar x\left(
       \theta_{\bar a}^{\delta_M}(\bar x)\to
       \theta_{\bar a}^{\delta_M+1}(\bar x)\right)
\]
is true in \(M\) and, among nonempty at most countable
structures, has exactly the models isomorphic to \(M\).
It can also be formed after naming any fixed finite
tuple, in which case it characterizes the pointed
isomorphism type, including repeated named coordinates.
It asserts uniqueness among countable models, not
among arbitrary uncountable models.
\end{hbtheorem}

\begin{hbcorollary}{thm:p03-morley-approximant-isomorphism}{proved=annex:p3.morley.approximant-isomorphism}{Isomorphism is the intersection of all countable approximants}
For at most countable \(M,N\),
\[
 M\cong N
       \quad\Longleftrightarrow\quad
 E_\alpha(M,N)\ \text{for every }\alpha<\omega_1.
\]
Each fixed \(E_\alpha\) is Borel on \(X_L\), but
the intersection here is indexed by all countable
ordinals, not by one countable list.
Closure of the Borel sets under countable intersections
does not make this an assertion that the whole
isomorphism relation is Borel.
\end{hbcorollary}

\subsection{Individual classes and a perfect family}
\label{sec:p03-morley-isomorphism-classes}

\begin{hbtheorem}{thm:p03-morley-isomorphism-complexity}{proved=annex:p3.morley.isomorphism-complexity}{Borel individual classes and analytic global isomorphism}
Each individual isomorphism class in \(X_L\) is Borel.
The global relation \(\cong\) on \(X_L\) is analytic:
it is the projection of a closed subset of
\(S_\infty\times X_L\times X_L\), where \(S_\infty\)
is the Polish group of permutations of \(\omega\).
Here analytic means projection of a Borel subset of
a product with a Polish space.
Borelness of each separate class does not by itself
prove Borelness of the global relation or a spectrum
dichotomy.
\end{hbtheorem}

\begin{hbtheorem}{thm:p03-morley-perfect-example}{proved=annex:p3.morley.perfect-example}{A complete perfect-set witness in an incomplete theory}
Let \(L=\{P_i:i<\omega\}\) consist of unary relations,
and let \(U\) say that every \(P_i\) is constant on
the domain:
\(\forall x\forall y(P_i(x)\leftrightarrow P_i(y))\).
For \(\eta\in2^\omega\), let \(M_\eta\) have domain
\(\omega\) and interpret \(P_i\) by all of \(\omega\)
when \(\eta(i)=1\), and by the empty set otherwise.
The map \(\eta\mapsto M_\eta\) is a homeomorphism
onto a nonempty perfect subset of \(X_L\) consisting
of pairwise nonisomorphic \(U\)-models.
It witnesses \(I(K_U)=2^{\aleph_0}\).
The theory \(U\) is incomplete: different models
decide \(\forall xP_i(x)\) differently.
The example is not a proof of the complete-theory
conjecture. The same binary patterns on any fixed
finite nonempty domain give continuum many finite
isomorphism types in this infinite language.
\end{hbtheorem}
